\chapter*{ИСПОЛЬЗОВАНИЕ СИСТЕМ ИСКУССТВЕННОГО ИНТЕЛЛЕКТА И ЛИЧНЫЙ ВКЛАД АВТОРА}
\addcontentsline{toc}{chapter}{Использование ИИ и личный вклад автора}
\sloppy

В~соответствии с~требованиями ИТМО к~ВКР настоящий раздел
описывает использование ИИ-инструментов обучающимся в~процессе подготовки работы
и~фиксирует личный вклад автора в~разработку системы.

\section*{Использование ИИ-инструментов}
\addcontentsline{toc}{section}{Использование ИИ-инструментов}

В~ходе подготовки работы автор использовал ИИ-инструменты как средства
поиска, чернового структурирования и~редакционной поддержки текста:

\begin{itemize}
  \item \textbf{OpenAI Codex / ChatGPT-5.4} --- для~структурирования
    отдельных разделов, поиска возможных формулировок и~редакционной
    шлифовки черновиков при~обязательной последующей ручной переработке;
  \item \textbf{Perplexity AI} --- для~поиска актуальных рыночных данных,
    проверки дат публикации и~навигации по~открытым источникам;
  \item \textbf{Claude Code} --- для ускорения работы по созданию технических спецификаций
    на разработку фронтенд и бекенд части системы.
\end{itemize}

Все ИИ-сгенерированные фрагменты проходили ручную проверку автора:
фактические утверждения перепроверялись по~первоисточникам,
а~финальные выводы и~исследовательские формулировки принимались
автором самостоятельно.

\textbf{ИИ-инструменты не~использовались} для~подмены первичных данных:
интервью, выводы исследования клиентов, экспериментальные артефакты, расчёты и~итоговые
исследовательские интерпретации были собраны, проверены и~оформлены
автором.

\section*{Личный вклад автора}
\addcontentsline{toc}{section}{Личный вклад автора}

Автор настоящей работы лично:

\begin{enumerate}
  \item \textbf{Провёл исследование предметной области}: собрал и
    проанализировал корпус из~45+ цитируемых академических
    и~индустриальных источников по~LLM, мультиагентным системам,
    RAG-подходам и~автоматизации создания образовательного контента;

  \item \textbf{Разработал архитектурное решение}: принял ключевые
    проектные решения по~выбору технологического стека (LangGraph,
    Dify/RAG, граф знаний, PostgreSQL-хранилище контрольных точек
    и результатов, Qwen-контур), спроектировал шестиузловой граф агентов
    и~выбранную архитектуру генерации с опорой на источник;

  \item \textbf{Разработал шаблоны запросов}: подготовил системные
    шаблоны запросов для~ADDIE, SAM и~Backward Design с~учётом
    требований таксономии Блума, особенностей русскоязычного
    ядрового корпуса и~зарубежного сравнительного контура;

  \item \textbf{Реализовал систему}: написал код LangGraph-графа агентов,
    шлюза LLM, FastAPI-приложения, слоя источниковой привязки
    и прослеживаемости, знака качества и снимка проверки дипломных метрик, настроил локальный
    и env-dev контуры проверки ИИ-методиста;

  \item \textbf{Провёл эксперименты}: сформировал ядровой корпус
    документов, подготовил 30 вопросно-ответных пар и~100 фактуальных утверждений,
    провёл измерения базовой конфигурации, парный B2/B4-прогон,
    детерминированную проверку фактов по схеме RAG и~сценарную
    оценку экономики единицы продукта;

  \item \textbf{Провёл Customer Development}: 20 глубинных B2B-интервью
    по~методу Mom Test и~систематизировал ключевые болевые точки рынка;

  \item \textbf{Подготовил доказательную базу ВКР}: оформил
    протокол внутренней и открытой апробации, машиночитаемые
    доказательные артефакты и автоматический контроль готовности, который
    проверяет согласованность текста диплома с кодом и результатами.
\end{enumerate}
